% GENERATED FILE. Edit canonical lessons, then run npm run generate:books.
% canonical-source-hash: fnv1a64:5c821c3b80db0050
% canonical-lessons: IT-C311-sposarsi, IT-C311-fede, IT-C311-velo, IT-C311-bouquet, IT-C311-bomboniera

\chapter{Il giorno del sì}
\label{ch:it-il-giorno-del-si}
\hlchaptermodality{311}

\begin{chapteropening}
\textbf{By the end of this chapter:} \emph{I can talk about a wedding day: getting married, the ring, the veil, the bouquet and the favours.}

Complete the last lesson of chapter 311: I can talk about a wedding day: getting married, the ring, the veil, the bouquet and the favours.
\end{chapteropening}

\section[sposarsi]{sposarsi — to get married}

\label{lesson:IT-C311-sposarsi}



\begin{quote}
Before the new one: say the Italian for to settle in, then the Italian for to return to one's country.
\end{quote}

\subsection*{You'll want to know: sposarsi}
\textbf{sposarsi} — \textquotedblleft{}to get married\textquotedblright{}.

\begin{quote}
Ci siamo sposati in un piccolo paese.
\end{quote}

\begin{quote}
\emph{We got married in a small village.}
\end{quote}

\emph{Sposare qualcuno} — to marry someone.

\begin{cousinweb}
\emph{sposarsi} is built on \emph{sposo} \textquotedblleft{}bridegroom\textquotedblright{}, inherited from Latin \emph{spōnsus} \textquotedblleft{}one promised\textquotedblright{}, from \emph{spondēre} \textquotedblleft{}to promise solemnly\textquotedblright{} (Lewis \& Short).
\end{cousinweb}

\begin{grammarlens}[title={What you've built}]
One more word for telling, explaining and describing at B1.
\end{grammarlens}

\subsection*{Your turn}
\begin{itemize}
\raggedright
  \item \emph{Say it:} \emph{sposarsi}
  \item \emph{Say it:} \emph{sposarsi}, once more
  \item \emph{Say it:} a sentence of your own with \emph{sposarsi}
  \item \emph{Recall:} say the Italian for to settle in, then the Italian for to return to one's country, then say \emph{sposarsi} again
\end{itemize}

\begin{tcolorbox}[breakable,colback=teal!4,colframe=teal!35!black,title={Before you move on}]
What does sposarsi mean? (\textquotedblleft{}To get married\textquotedblright{}.) Say it once more.
\end{tcolorbox}
\section[la fede]{la fede — the wedding ring; faith}

\label{lesson:IT-C311-fede}



\begin{quote}
Before the new one: say the Italian for to return to one's country, then the Italian for to get married.
\end{quote}

\subsection*{You'll want to know: la fede}
\textbf{la fede} — \textquotedblleft{}the wedding ring; faith\textquotedblright{}.

\begin{quote}
Porta la fede alla mano sinistra.
\end{quote}

\begin{quote}
\emph{She wears her wedding ring on her left hand.}
\end{quote}

The ring is a sign of \emph{fede}, faith and trust.

\begin{cousinweb}
\emph{fede} is inherited from Latin \emph{fidēs} \textquotedblleft{}trust, faith\textquotedblright{} (EtymDB): the ring is a pledge of faith.
\end{cousinweb}

\begin{grammarlens}[title={What you've built}]
One more word for telling, explaining and describing at B1.
\end{grammarlens}

\subsection*{Your turn}
\begin{itemize}
\raggedright
  \item \emph{Say it:} \emph{la fede}
  \item \emph{Say it:} \emph{la fede}, once more
  \item \emph{Say it:} a sentence of your own with \emph{la fede}
  \item \emph{Recall:} say the Italian for to return to one's country, then the Italian for to get married, then say \emph{la fede} again
\end{itemize}

\begin{tcolorbox}[breakable,colback=teal!4,colframe=teal!35!black,title={Before you move on}]
What does la fede mean? (\textquotedblleft{}The wedding ring; faith\textquotedblright{}.) Say it once more.
\end{tcolorbox}
\section[il velo]{il velo — the veil}

\label{lesson:IT-C311-velo}



\begin{quote}
Before the new one: say the Italian for to get married, then the Italian for the wedding ring.
\end{quote}

\subsection*{You'll want to know: il velo}
\textbf{il velo} — \textquotedblleft{}the veil\textquotedblright{}.

\begin{quote}
La sposa ha un velo lunghissimo.
\end{quote}

\begin{quote}
\emph{The bride has a very long veil.}
\end{quote}

\emph{Il velo da sposa} — the bridal veil.

\begin{cousinweb}
\emph{velo} is inherited from Latin \emph{vēlum} \textquotedblleft{}a sail, a curtain\textquotedblright{} (EtymDB).
\end{cousinweb}

\begin{grammarlens}[title={What you've built}]
One more word for telling, explaining and describing at B1.
\end{grammarlens}

\subsection*{Your turn}
\begin{itemize}
\raggedright
  \item \emph{Say it:} \emph{il velo}
  \item \emph{Say it:} \emph{il velo}, once more
  \item \emph{Say it:} a sentence of your own with \emph{il velo}
  \item \emph{Recall:} say the Italian for to get married, then the Italian for the wedding ring, then say \emph{il velo} again
\end{itemize}

\begin{tcolorbox}[breakable,colback=teal!4,colframe=teal!35!black,title={Before you move on}]
What does il velo mean? (\textquotedblleft{}The veil\textquotedblright{}.) Say it once more.
\end{tcolorbox}
\section[il bouquet]{il bouquet — the bouquet}

\label{lesson:IT-C311-bouquet}



\begin{quote}
Before the new one: say the Italian for the wedding ring, then the Italian for the veil.
\end{quote}

\subsection*{You'll want to know: il bouquet}
\textbf{il bouquet} — \textquotedblleft{}the bouquet\textquotedblright{}.

\begin{quote}
Alla fine la sposa lancia il bouquet.
\end{quote}

\begin{quote}
\emph{At the end the bride throws the bouquet.}
\end{quote}

From French; said \emph{bu-KEH}.

\begin{cousinweb}
\emph{bouquet} was borrowed from French \emph{bouquet} \textquotedblleft{}a little wood, a bunch\textquotedblright{}; EtymDB traces the French word to Old French \emph{bochet}, a little wood, and has no Italian entry.
\end{cousinweb}

\begin{grammarlens}[title={What you've built}]
One more word for telling, explaining and describing at B1.
\end{grammarlens}

\subsection*{Your turn}
\begin{itemize}
\raggedright
  \item \emph{Say it:} \emph{il bouquet}
  \item \emph{Say it:} \emph{il bouquet}, once more
  \item \emph{Say it:} a sentence of your own with \emph{il bouquet}
  \item \emph{Recall:} say the Italian for the wedding ring, then the Italian for the veil, then say \emph{il bouquet} again
\end{itemize}

\begin{tcolorbox}[breakable,colback=teal!4,colframe=teal!35!black,title={Before you move on}]
What does il bouquet mean? (\textquotedblleft{}The bouquet\textquotedblright{}.) Say it once more.
\end{tcolorbox}
\section[la bomboniera]{la bomboniera — the wedding favour}

\label{lesson:IT-C311-bomboniera}



\begin{quote}
Before the new one: say the Italian for the veil, then the Italian for the bouquet.
\end{quote}

\subsection*{You'll want to know: la bomboniera}
\textbf{la bomboniera} — \textquotedblleft{}the wedding favour\textquotedblright{}.

\begin{quote}
Abbiamo scelto bomboniere di vetro.
\end{quote}

\begin{quote}
\emph{We chose glass wedding favours.}
\end{quote}

A small gift with sugared almonds, given at weddings and christenings.

\begin{cousinweb}
\emph{bomboniera} is usually traced to French \emph{bonbonnière} \textquotedblleft{}a sweet box\textquotedblright{}, built on \emph{bonbon} \textquotedblleft{}sweet\textquotedblright{}, from \emph{bon} \textquotedblleft{}good\textquotedblright{}, Latin \emph{bonus} (Lewis \& Short).
\end{cousinweb}

\begin{grammarlens}[title={What you've built}]
One more word for telling, explaining and describing at B1.
\end{grammarlens}

\subsection*{Your turn}
\begin{itemize}
\raggedright
  \item \emph{Say it:} \emph{la bomboniera}
  \item \emph{Say it:} \emph{la bomboniera}, once more
  \item \emph{Say it:} a sentence of your own with \emph{la bomboniera}
  \item \emph{Recall:} say the Italian for the veil, then the Italian for the bouquet, then say \emph{la bomboniera} again
\end{itemize}

\begin{tcolorbox}[breakable,colback=teal!4,colframe=teal!35!black,title={Before you move on}]
What does la bomboniera mean? (\textquotedblleft{}The wedding favour\textquotedblright{}.) Say it once more.
\end{tcolorbox}
